\chapter{Direitos do estudante}
\DoPToC
\section{Mobilidade acadêmica}

    \begin{wrapfigure}{L}{0.5\textwidth}
        \centering
        \includegraphics[width=0.49\textwidth]{mobi}
    \end{wrapfigure}

 O Programa de Mobilidade Acadêmica foi criado com o objetivo de permitir que estudantes de graduação de Instituições Federais de Ensino Superior (IFES) \index{IFES} brasileiras pudessem realizar intercâmbio entre elas. Para garantir sua participação o estudante precisa ter cursado todas as disciplinas do primeiro ano, ou seja, 1$^{\mbox{o}}$ e 2$^{\mbox{o}}$ semestres letivos do curso na Instituição de origem, e possuir no máximo uma (01) reprovação por período letivo.

O processo de mobilidade acadêmica não se aplica a pedidos de transferência de alunos entre as Instituições Federais de Ensino Superior, que normalmente é feito por Concurso de Vagas Residuais.
    \begin{enumerate}
        \item O vínculo é temporário com a Instituição receptora, dependendo da existência de disponibilidade de vaga e das possibilidades de inscrição na(s) disciplina(s) pretendida(s);
        \item O prazo máximo para afastamento é de 1 ano letivo, e em caráter excepcional, a critério da Instituição receptora, poderá haver renovação, sucessiva ou intercalada, do vínculo temporário, por até mais um período letivo.
    \end{enumerate}

Os estudantes interessados em participar, deverão consultar o Coordenador do seu respectivo Curso para elaborar o plano de estudos, que envolve escolhas de disciplinas pretendidas sendo correlacionadas com as ementas de sua IFES de origem, para que ocorra o aproveitamento quando o estudante retornar a sua IFES de origem.
   
O próximo passo é o preenchimento do requerimento (formulário concedido pela SGC) e em anexo especificado as disciplinas que serão aproveitadas e sua carga horária, entregar ao Colegiado para posteriormente o próprio colegiado encaminhar para a PROGRAD.

A Pró-Reitoria de Graduação irá consulta o Pró-Reitor da IFES pleiteada caso a IFES aceitar receber o estudante, após a resposta da IFES a Pró-Reitoria de Graduação comunica ao colegiado para informar ao estudante interessado.

 De acordo com o convênio, a UFBA não se responsabilizará pelas despesas dos participantes no programa.
 
 Os prazos obedecem ao disposto no Calendário Acadêmico, seja para alunos da UFBA ou de outras IFES.
 
 \subsection{Mobilidade Acadêmica Internacional}
 
 A Mobilidade Acadêmica Internacional acontece através de parcerias estabelecidas entre Universidades Internacionais e Programas Brasileiros de Intercâmbio relacionados com Universidades brasileiras.

 
\section{Aproveitamento de disciplina}

\begin{wrapfigure}{l}{0.2\textwidth}
        \centering
        \includegraphics[width=0.19\textwidth]{aproveite}
    \end{wrapfigure}

A dispensa de cursar, para efeito de integralização do currículo, determinado componente curricular ocorrerá mediante aproveitamento de estudos realizados nesta ou em outra instituição de ensino superior, ou através de avaliação de conhecimento prévio.
Para que o aproveitamento de disciplinas seja concedido, o estudante deverá:
\begin{itemize}
    \item Apresentar requerimento ao Colegiado do curso, instruído com o seu histórico escolar, ementa(s) e conteúdo(s) programático(s) do(s) estudo(s) realizado(s);
    \item No caso de documentos oriundos de instituições estrangeiras, estes deverão estar acompanhados pelas respectivas traduções juramentadas e ser autenticados pela autoridade consular brasileira no país de origem, exceto em caso de programa de intercâmbio, acordo ou convênio com a UFBA, a critério do Colegiado do curso.
\end{itemize}

No caso de dispensa de componente curricular obrigatório, o conteúdo programático e carga horária do componente cursado deverão corresponder ao mínimo de 75\% (setenta e cinco por cento) daqueles referentes ao componente dispensado, exceto em situações especiais, a critério do Colegiado.

No caso de dispensa de carga horária/crédito em componentes curriculares optativos ou livres, o Colegiado definirá o total a ser dispensado.

\section{Exame de suficiência}

O estudante pode também solicitar o aproveitamento de disciplina através de um exame de conhecimento prévio. Para isto, o estudante deverá procurar o Colegiado do curso.

 \begin{wrapfigure}{r}{0.4\textwidth}
        \centering
        \includegraphics[width=0.39\textwidth]{prova}
    \end{wrapfigure}

O exame para a avaliação do conhecimento prévio deverá ser solicitado pelo Colegiado do curso ao órgão responsável pela oferta do componente curricular que, por sua vez, deverá constituir uma Comissão de Avaliação, composta por, no mínimo, três (03) professores, sendo um deles, obrigatoriamente, docente responsável pelo componente curricular no qual o aluno será avaliado. 

O exame de conhecimento prévio deverá considerar como referência a ementa do componente curricular, adotando processos avaliativos compatíveis com as características do componente a ser dispensado.

\section{Trancamento de matrícula}

    \begin{wrapfigure}{r}{0.4\textwidth}
        \centering
        \includegraphics[width=0.39\textwidth]{tialara}
    \end{wrapfigure}

O trancamento de matrícula ou o trancamento total de inscrição em componentes curriculares, na graduação, não poderá ultrapassar 50\% (cinquenta por cento) do tempo máximo do fluxograma mais atualizado do curso. 
    
Se dentro do prazo não é necessário uma justificativa apenas o preenchimento do formulário, caso contrário, você deve explicar o motivo pelo qual você deseja interromper seus estudos
    
O trancamento parcial de inscrição em componentes curriculares, na graduação, só será permitido até o máximo de 20\% (vinte por cento) da carga horária total do respectivo curso.

O trancamento de matrícula, por motivo de saúde, deverá ser comprovado por atestado médico. Esse trancamento poderá ser solicitado a qualquer tempo e não será computado na integralização do curso. Além disso, tanto o prazo de duração do referido trancamento, quanto o retorno do estudante às atividades acadêmicas deverão considerar informações constantes em atestado médico.
    
O trancamento de matrícula, exceto por motivo de saúde, poderá ser interrompido a qualquer época, a pedido do aluno, para inscrição em componentes curriculares no semestre letivo seguinte à interrupção.

\section{Afastamento}

O afastamento, poderá ser solicitado pelo aluno e será concedido a este, quando requerido dentro do período estabelecido na agenda acadêmica. Este afastamento não poderá ultrapassar 50\% (cinquenta por cento) do tempo máximo do fluxograma mais atualizado curso. Ou seja, se o curso tiver um fluxograma com duração de quatro anos, o trancamento máximo da matrícula deverá ser de dois anos.

O aluno será notificado pelo sistema informatizado de registro e controle acadêmico quanto à sua situação referente ao percentual de trancamento, quando o mesmo alcançar setenta por cento do tempo limite definido.

\begin{itemize}
    \item Por motivos de saúde, comprovados por atestado médico, o trancamento de matrícula poderá ser solicitado a qualquer tempo e não será computado na integralização do curso. 
    \item Tanto o prazo de duração do referido trancamento, quanto o retorno do estudante às atividades acadêmicas deverão considerar informações constantes em atestado médico.
\end{itemize}
O afastamento, exceto por motivo de saúde, poderá ser interrompido a qualquer época, a pedido do aluno, para inscrição em componentes curriculares no semestre letivo seguinte à interrupção. 

\section{Regimes excepcionais}

Será concedido regime excepcional ao estudante que se enquadrar nas determinações do Decreto-Lei 1.044/69, da Lei n$^{\mbox{o}}$ 6.202/75 e nas normas estabelecidas pela Resolução 09/2009/CEPE.

\begin{wrapfigure}{r}{0.3\textwidth}
        \centering
        \includegraphics[width=0.29\textwidth]{gravida}
    \end{wrapfigure}
São considerados merecedores de tratamento excepcional os alunos de qualquer nível de ensino, portadores de afecções congênitas ou adquiridas, infecções, traumatismo ou outras condições mórbitas, determinando distúrbios agudos ou agudizados, caracterizados por: 

    \begin{enumerate}
    
        \item Incapacidade física relativa, incompatível com a frequência aos trabalhos escolares, desde que se verifique a conservação das condições intelectuais e emocionais necessárias para o prosseguimento da atividade escolar em novos moldes;
        \item Ocorrência isolada ou esporádica;
        \item Duração que não ultrapasse o máximo ainda admissível, em cada caso, para a continuidade do processo pedagógico de aprendizado, atendendo a que tais características se verificam, entre outros, em casos de síndromes hemorrágicos (tais como a hemofilia), asma, cartide, pericardites, afecções osteoarticulares submetidas a correções ortopédicas, nefropatias agudas ou subagudas, afecções reumáticas, entre outras.
    \end{enumerate}

\begin{itemize}
    \item A esses estudantes deverão serem atribuídos como compensação da ausência às aulas, exercício domiciliares com acompanhamento da instituição, sempre que compatíveis com o seu estado de saúde e as possibilidades do estabelecimento.
    \item O regime de exceção dependerá de laudo médico elaborado por autoridade oficial do sistema educacional. 
    \item Em caso de gravidez, de acordo com a Lei 6.202/75 a estudante,
a partir do oitavo mês de gestação e
durante 3(três) meses, ficará assistida pelo regime
de exercício domiciliar instituído pelo Decreto-
Lei 1.044/69. 
    \item Será da competência do Diretor do estabelecimento a autorização, à autoridade superior imediata, do regime de exceção.
    \item Serão de responsabilidade do estudante o acompanhamento da matéria ministrada e o cumprimento das atividades planejadas e de outras obrigações inerentes, durante o período de excepcionalidade. 
\end{itemize}

\section{Dilatação de prazo}

A dilatação do prazo consiste em estender o período de curso que por algum motivo não fora concluído em seu tempo máximo. O estudante poderá requerer dilação do prazo máximo para conclusão de seu curso a partir do momento que apresentar motivos que possam ser comprovados, como por exemplo, motivo de saúde mediante a atestado ou se o estudante estava em intercâmbio. Esse requerimento só poderá ser solicitado se o estudante se encontrar em seu último período e é concedida a dilatação do prazo correspondente a 50\% (cinquenta por cento) do tempo de curso. A solicitação poderá ser feita na Secretaria Geral dos Cursos através de formulário e documento de justificativa da dilatação, desde que dentro do período estipulado no Calendário Acadêmico.

\section{Orientação Acadêmica}

A orientação acadêmica é uma relação estabelecida entre professores e estudantes que consiste no encaminhamento  de um grupo de estudantes para um professor orientador, durante o seu tempo de curso na intenção de ajudá-los em questões acadêmicas, profissionais e administrativas. A orientação acadêmica colabora com os estudantes nas suas futuras escolhas, que podem ser profissionais ou acadêmicas, através de discussões sobre as disciplinas a serem cursadas, o mercado de trabalho, as áreas de atuação e pesquisas. O professor orientador poderá oferecer também sugestões para atividades de extensão, sanar dúvidas sobre o funcionamento do departamento, e eliminar outras dúvidas que possam surgir.